\section{Control Rights and the Commitment Problem}\label{sec:control}

The static model takes Europe's ownership of the chokepoint as given. The scenario's 2031 crisis is about ownership itself. Washington proposes to fold ASML into a joint holding company in which it holds a controlling vote over output, customers, and technology transfer. This section asks when a patron prefers to buy or coerce control of a chokepoint rather than pay rents to its holder, and what that does to the holder.

The answer turns on a commitment problem of the kind \citet{fearon1995rationalist} and \citet{powell2006war} identify in shifting-power settings, applied to an object whose ownership shapes future bargaining power \citep{fearon1996bargaining}. Europe's exit premium can change over time. Rearmament lowers $\kappa S$; Chinese credit lines and market offers raise $\omega$; rationing of loyal allies, if it comes, removes the conditional part of the AI lever (\cref{prop:statics}b). If Europe's premium is expected to turn positive, Washington faces future rents, and Europe cannot credibly promise to forgo them. Control rights remove the problem. Ownership determines who gets to choose alignment later, and the patron can buy it while its own threats still bite.

\subsection{Setup}

There are two periods, $t \in \{1,2\}$, and a discount factor $\delta \in (0,1]$. In each period in which Europe holds control, the static game of \cref{sec:model} is played with period-specific parameters, giving exit premium $Q_t$ and AI exposure $\alpha_t$. We hold the AI lever conditional in both periods ($\chi_{ai,t}=1$, so $\nu_t < k_A$) and let $Q_t$ move for the other reasons above. We also assume that the United States wins each static contest even after a cut ($V_t > Q_t + \alpha_t$ and $B_U \geq V_t$).

At the start of period 1 the United States may issue an ultimatum: Europe transfers control over ASML's customers, output, and technology for both periods in exchange for a payment $T \geq 0$ made at signing; if Europe refuses, its access to frontier AI is cut for all of period 1. The United States chooses how public to make the ultimatum, which fixes the audience cost $m \in [0,\bar m]$ it would pay for not carrying out the threat \citep{fearon1994domestic}. If Europe accepts, it bears a sovereignty cost $\sigma \geq 0$, which includes the domestic political cost of surrendering a national champion and any capitalized loss of security from relocating capacity to US soil, and China punishes the transfer with $r^c \geq 0$, the part of its threatened punishment that passes the test of \cref{lem:credible}(iv). Once control has passed, the United States has no reason to cut Europe off: its ex post gain from doing so is negative, so the assurance half of the ultimatum is credible \citep{cebul2021coercion}. If Europe refuses, the United States decides whether to cut. Cutting a Europe that remains aligned yields the direct ex post gain $g_N \equiv \nu_1 - k_A$, which is negative because $\chi_{ai,1}=1$ requires $\nu_1<k_A$. It also removes the AI term from Europe's period-1 exit costs, so the period-1 exit premium rises from $Q_1$ to $Q_1 + \alpha_1$. We assume access is restored in period 2; \cref{lem:selflimit} shows that, when $Q_2 \geq 0$, a longer cut would cost Washington rent without adding to its bite. Finally, the United States issues an ultimatum only if doing so strictly raises its payoff, the limit of an arbitrarily small cost of issuing one.

Payoffs, normalized as before, are:
\begin{center}
\begin{tabular}{lll}
\toprule
Outcome & Europe & United States\\
\midrule
No ultimatum & $[Q_1]^+ + \delta[Q_2]^+$ & $-[Q_1]^+ - \delta[Q_2]^+$\\
Accept & $T - \sigma - r^c$ & $-T$\\
Refuse, cut & $[Q_1+\alpha_1]^+ - \alpha_1 + \delta[Q_2]^+$ & $g_N - [Q_1+\alpha_1]^+ - \delta[Q_2]^+$\\
Refuse, no cut & $[Q_1]^+ + \delta[Q_2]^+$ & $-m - [Q_1]^+ - \delta[Q_2]^+$\\
\bottomrule
\end{tabular}
\end{center}
Every row omits China's per-period punishment of Atlantic alignment, $\chi_{r,t} r_t$, which Europe bears whoever controls ASML; $r^c$ is China's additional punishment of the transfer itself.

\subsection{The bite of a credible threat}

\begin{lemma}[Credible ultimatum and effective bite]\label{lem:bite}
After a refusal, the United States cuts Europe off if and only if
\begin{equation}\label{eq:credcut}
m + g_N \;\geq\; [Q_1+\alpha_1]^+ - [Q_1]^+ .
\end{equation}
When it does, Europe's net loss is the \emph{effective bite}
\begin{equation}\label{eq:bite}
b_1 \equiv \alpha_1 - \big([Q_1+\alpha_1]^+ - [Q_1]^+\big) = \min\{\alpha_1, [-Q_1]^+\}.
\end{equation}
\end{lemma}

\Cref{lem:bite} applies the test of \cref{lem:credible} to a sanction imposed before a contest rather than after one, so the sanction's effect on the rent the United States must later pay enters the comparison. Two features matter. First, without an audience cost the threat is not credible: with $m=0$ and $g_N<0$, cutting off a Europe that has merely refused to sell control costs Washington something and gains it nothing. That is why the scenario's White House has to go public. Second, a punishment hurts its target only up to the target's slack. Cutting AI access removes one of the reasons Europe stays aligned. Once $Q_1 + \alpha_1>0$, the rest of the cut comes back to Europe as rent, because Washington must now pay to keep it from exiting. A large stick adds nothing beyond the point where using it would make exit attractive, which is the logic of marginal deterrence \citep{stigler1970optimum}.

\begin{lemma}[Self-limiting sanctions]\label{lem:selflimit}
Extending the cut into period 2 costs the United States $[Q_2+\alpha_2]^+ - [Q_2]^+$ in additional rent, which equals $\alpha_2$ when $Q_2 \geq 0$. When $Q_2 \geq 0$ the extension imposes no additional net loss on Europe. A permanent cut is credible only if the period-2 audience cost is at least this rent less the period-2 gain $\nu_2 - k_A$.
\end{lemma}

The scenario's threat that Europe will ``lose all access, present and future, to American AI'' is therefore not credible inside the model unless Washington has staked a great deal of reputation on it. A permanent cut hands Europe exactly the exit option the cut was meant to deter, at the moment Europe's future premium is positive.

\subsection{Preventive control}

\begin{proposition}[Preventive control]\label{prop:control}
Under complete information, the United States issues a public ultimatum, and Europe accepts, if and only if
\begin{enumerate}[label=(\roman*)]
\item Europe's future exit premium is positive, $Q_2>0$;
\item the threat can be made credible, $\bar m + g_N \geq [Q_1+\alpha_1]^+ - [Q_1]^+$; and
\item the effective bite exceeds the deadweight cost of transferring control, $b_1 > \sigma + r^c$.
\end{enumerate}
The payment is $T^* = [\delta Q_2 + \sigma + r^c - b_1]^+$ and Europe's payoff is $\max\{\delta Q_2 - b_1,\, -\sigma - r^c\}$.
\end{proposition}

Because $b_1 \leq [-Q_1]^+$, condition (iii) requires $Q_1 < -(\sigma + r^c)$. In the two-period model, Washington must act while Europe's exit premium is still well below zero, and it acts because that premium is expected to become positive. A patron that faces no future rents gains nothing from control, so any cost of demanding it, however small, deters the demand. The tie-breaking assumption matters on an open set of parameters: when $Q_2\leq 0$ and $b_1\geq\sigma+r^c$, Europe would accept a demand for control without payment, so a patron that placed any value on control for its own sake would demand it even without future rents. The same logic extends to many periods.

\begin{corollary}[Windows]\label{cor:windows}
Consider an $N$-period version in which, at any date $t$ at which Europe still holds control, the United States may issue an ultimatum that transfers control for all remaining dates; a refusal is met by cutting AI access at date $t$ only; the bite is $b_t \equiv \min\{\alpha_t, [-Q_t]^+\}$; and the United States issues only when it strictly gains. It issues an ultimatum at some date if and only if there is a date $t$ with $\bar m + \nu_t - k_A \geq [Q_t+\alpha_t]^+-[Q_t]^+$ and $b_t>\sigma+r^c$, and a later date $t'>t$ with $Q_{t'}>0$.
\end{corollary}

\begin{figure}[t]
\centering
\includegraphics[width=0.95\textwidth]{figures/fig_windows.pdf}
\caption{Windows for preventive control (illustrative units). Shaded blue: dates at which a credible US threat has a bite larger than $\sigma + r^c$, which requires $Q_t < -(\sigma+r^c)$. Shaded orange: Europe's window, $Q_t>0$. In panel (a) a US window precedes Europe's, so Washington demands control before Europe's window opens (\cref{cor:windows}). In panel (b) Europe's premium rises but never turns positive, there are no future rents to remove, and no demand is made.}
\label{fig:windows}
\end{figure}

At dates in the US window, where the cut is credible and the bite exceeds $\sigma + r^c$, the exit premium is below $-(\sigma+r^c)$; at dates in Europe's window it is positive. The demand need not come at the end of the US window, or even inside it. Because a refusal also exposes Europe to a later demand, Washington can find it cheaper to act earlier, even at a date when the bite alone would not justify an ultimatum. Under complete information the speed of the shift therefore does not affect whether Washington acts, only when it acts and what it pays. Two features of \cref{prop:control} cut against intuition. China's threat to punish the transfer, $r^c$, protects European ownership, because it raises the bite Washington needs; this requires that China would actually carry out the punishment once the transfer is done. And Europe's sovereignty costs, which commentators treat as an obstacle to a deal, play the same role: a Netherlands that would pay a high domestic price for surrendering ASML is harder to coerce into doing so.

\subsection{Private sovereignty costs and breakdown}

The scenario's endgame is a standoff, not a signature. Under complete information the ultimatum is either accepted or never issued. Breakdown requires that Washington be uncertain about what surrender would cost the European governments involved, which is plausible given coalition politics in The Hague and Berlin. Let $\sigma$ be Europe's private information, uniformly distributed on $[0,\bar\sigma]$ with $\bar\sigma>0$, and suppose conditions (i) and (ii) of \cref{prop:control} hold. Write $R \equiv \delta Q_2$, $\tilde b \equiv b_1 - r^c > 0$, and $e \equiv -g_N + [Q_1 + \alpha_1]^+ > 0$ for what it costs Washington to carry out the cut.

\begin{proposition}[Breakdown]\label{prop:breakdown}
If the United States issues the ultimatum, Europe accepts if and only if $\sigma \leq \hat\sigma$, where
\begin{equation}\label{eq:sstar}
\hat\sigma = \max\Big\{\tilde b - R,\; \min\big\{\bar\sigma,\, \tfrac{\tilde b+e}{2}\big\}\Big\}.
\end{equation}
Let $F^* \equiv \min\{\hat\sigma/\bar\sigma, 1\}$. The United States issues the ultimatum if and only if $F^*(\tilde b - \hat\sigma) > (1-F^*)e$. If it does, Europe refuses with probability $P_b = 1 - F^*$, the cut is carried out, and Europe keeps ASML. Conditional on the ultimatum being issued, $P_b$ is nonincreasing in the bite $b_1$ and in the execution cost $e$ (each holding the other fixed), and nondecreasing in China's punishment of the transfer $r^c$ and in $\bar\sigma$. When $\min\{\bar\sigma, (\tilde b+e)/2\} < \tilde b - R < \bar\sigma$, the United States offers nothing and $P_b$ is strictly increasing in Europe's future premium $R$; when $\tilde b - R \geq \bar\sigma$, it offers nothing and every type accepts.
\end{proposition}

Breakdown is the scenario's ``neither'' outcome: Europe refuses both great powers, loses frontier access, and keeps the chokepoint. Once an ultimatum is issued, refusal is more likely when China credibly threatens to punish a transfer, when the European domestic cost of surrender may be high, and when Washington's threat is cheap to execute. The last point is the usual risk-return trade-off \citep{fearon1995rationalist, powell1999shadow}: a patron that loses little from a standoff makes tougher offers and is refused more often. A stronger Chinese threat, or a higher ceiling on Europe's sovereignty costs, can also deter the ultimatum altogether.

\subsection{Promises or threats}

Europe could try to remove the commitment problem itself. An allied export-control compact, dual-key licensing that gives Washington a veto over China-bound sales, or treaty-entrenched EU controls would bind Europe to Atlantic alignment in period 2 without transferring ownership. Suppose Europe can adopt such a compact unilaterally before period 1, at cost $\sigma^b$, and China punishes it with $r^b$, again the credible part of its threat.

\begin{remark}[Self-binding]\label{rem:bind}
When the conditions of \cref{prop:control} hold, Europe strictly prefers to bind itself if and only if $\sigma^b + r^b < \sigma + r^c$ and $\sigma^b + r^b < b_1 - \delta Q_2$. Binding removes the US motive for an ultimatum and keeps ownership in Europe, but it forfeits the future rents that a positive $Q_2$ would have earned. If control had any value to Washington in itself, binding would not prevent a demand for it.
\end{remark}

A credible promise not to exit and a credible exit threat cannot coexist: the compact works by setting Europe's future exit premium to minus infinity. It does not remove the other source of leverage, the threat to withhold servicing from the United States itself, which \cref{sec:holdup} analyzes. Binding against exit also removes the preventive part of Washington's stake; we conjecture, but do not show, that this could make the United States more willing to concede a hold-up dispute.
